\chapter{Food Intolerance}

\section*{Definition}
Food intolerance is an adverse reaction to a food or food component that does not involve the immune system, distinguished from food allergy by the absence of IgE-mediated mechanisms. Symptoms typically are dose-dependent, delayed in onset (hours to days), and reproducible with the offending food.

\section*{What Happens in Your Body}
Deficiency of digestive enzymes (e.g., lactase deficiency in lactose intolerance), pharmacologic effects of food components (e.g., vasoactive amines in histamine intolerance), or direct mucosal irritation (e.g., capsaicin) triggers gastrointestinal and systemic symptoms. Altered gut microbiota and impaired intestinal barrier function may contribute.

\section*{Causes}
\begin{itemize}
  \item \textbf{Enzymatic:} Lactose intolerance (lactase deficiency), sucrase-isomaltase deficiency, pancreatic exocrine insufficiency
  \item \textbf{Pharmacologic:} Caffeine (jitters, palpitations), tyramine (migraine trigger in MAO-I users), histamine (scombroid poisoning, fermented foods), sulfites (asthma exacerbation), MSG (transient symptoms in sensitive individuals)
  \item \textbf{FODMAPs:} Fermentable oligo-, di-, mono-saccharides and polyols (fructose, lactose, fructans, galactans, polyols) causing gas and bloating in IBS people
  \item \textbf{Other:} Gluten sensitivity (non-celiac gluten sensitivity), artificial sweeteners (sorbitol, xylitol), food colorings and preservatives, nickel allergy (systemic contact dermatitis)
\end{itemize}

\section*{When to See a Doctor}
See your doctor if:
\begin{itemize}
  \item Severe abdominal pain, vomiting, diarrhea with dehydration, or weight loss
  \item Blood in stool, iron-deficiency anemia, or nutritional deficiencies
  \item Symptoms starting in infancy or childhood with poor growth
  \item Suspected celiac disease (warrants serology before gluten elimination)
  \item Food-triggered anaphylaxis or urticaria suggests IgE-mediated allergy, not intolerance
\end{itemize}

\section*{Related Symptoms}
\begin{itemize}
  \item Bloating, gas, abdominal pain, diarrhea, nausea, vomiting, headache, fatigue, brain fog, skin rash
\end{itemize}
\textbf{Important:} Unlike food allergies, food intolerances are not life-threatening and do not involve IgE antibodies; skin prick tests and serum-specific IgE tests are negative.

\section*{Self-Care / Home Management}
\begin{itemize}
  \item Keep a detailed food and symptom diary to identify trigger foods and dose thresholds
  \item Trial elimination diet: remove suspected foods for 2--4 weeks, then reintroduce systematically
  \item Lactose intolerance: use lactase enzyme supplements (drops or tablets) or choose lactose-free dairy products
  \item Low-FODMAP diet under dietician guidance for suspected fermentable carbohydrate intolerance
  \item Read food labels for hidden sources of sulfites, MSG, artificial sweeteners, or other triggers
\end{itemize}

\section*{How Your Doctor Will Evaluate You}
\begin{itemize}
  \item Clinical history and food diary most valuable; identify temporal pattern between food intake and symptoms
  \item Lactose hydrogen breath test for suspected lactose intolerance (elevated H2 after 25--50 g lactose load)
  \item Celiac serology (tTG-IgA, total IgA) to rule out celiac disease before gluten elimination
  \item Fructose or lactulose breath testing for FODMAP malabsorption
  \item Elimination diet and double-blind placebo-controlled food challenge (gold standard, rarely performed)
\end{itemize}

\section*{Treatment Options}
\begin{itemize}
  \item Dietary avoidance of identified trigger foods, tailored to individual tolerance thresholds
  \item Enzyme replacement: lactase for lactose intolerance, alpha-galactosidase (Beano) for gas from beans/vegetables
  \item Probiotics may improve symptoms in lactose intolerance and functional bloating through microbiome modulation
  \item Low-FODMAP diet under dietician supervision for IBS-related food intolerance
  \item Antispasmodics or peppermint oil for symptom flares; no pharmacotherapy eliminates the underlying intolerance
\end{itemize}

\section*{References}
\begin{thebibliography}{99}
\bibitem{ref1} Lomer MCE. "Review Article: The Aetiology, Diagnosis, Mechanisms and Clinical Evidence for Food Intolerance." Aliment Pharmacol Ther. 2015;41(3):262--275.
\bibitem{ref2} Turnbull JL, Adams HN, Gorard DA. "Review Article: The Diagnosis and Treatment of Food Allergy and Food Intolerance." Aliment Pharmacol Ther. 2015;41(1):3--25.
\bibitem{ref3} Boyce JA, Assa'ad A, Burks AW, et al. "Guidelines for the Diagnosis and Management of Food Allergy in the United States." J Allergy Clin Immunol. 2010;126(6 Suppl):S1--S58.
\bibitem{ref4} Ortolani C, Pastorello EA. "Food Allergies and Food Intolerances." Best Pract Res Clin Gastroenterol. 2006;20(3):467--483.
\bibitem{ref5} Zopf Y, Baenkler HW, Silbermann A, et al. "The Differential Diagnosis of Food Intolerance." Dtsch Arztebl Int. 2009;106(21):359--370.
\end{thebibliography}
\clearpage
